% Generated draft: requires mathematical and source review.
\documentclass[11pt]{article}
\usepackage[margin=25mm]{geometry}
\usepackage{fontspec}
\setmainfont{Latin Modern Roman}
\usepackage{amsmath,amssymb,mathtools,unicode-math}
\setmathfont{Latin Modern Math}
\usepackage{enumitem,needspace,xurl,hyperref}
\setlength{\parindent}{0pt}
\setlength{\parskip}{5pt}
\setlength{\emergencystretch}{3em}
\newcommand{\N}{\mathbb{N}}
\newcommand{\Z}{\mathbb{Z}}
\newcommand{\Q}{\mathbb{Q}}
\newcommand{\R}{\mathbb{R}}
\newcommand{\C}{\mathbb{C}}
\newcommand{\F}{\mathbb{F}}
\newcommand{\A}{\mathbb{A}}
\newcommand{\PP}{\mathbb{P}}
\newcommand{\E}{\mathbb{E}}
\newcommand{\eps}{\varepsilon}
\newcommand{\dd}{\,\mathrm{d}}
\begin{document}
\section*{UnsolvedMath ID: 20002468}
\textbf{Source problem code:} catalogue entry
\subsection{Definitions and mathematical statement}
Patrick Bernard's Question 5 concerns an analytic Lagrangian $L:\mathbb T^n\times\R^n\times\mathbb T\to\R$, $L=L(x,v,t)$, satisfying Mather's conditions. In the conventional time-periodic Tonelli interpretation, $\mathbb T=\R/\Z$, $D^2_{vv}L>0$, $L$ is superlinear in $v$, and the Euler--Lagrange flow has the completeness needed for the minimizing-measure theory. ``Positive definite'' concerns the velocity Hessian, rather than the sign of $L$. Let $\mathcal I$ be the invariant probability measures of the extended flow with finite relevant action, and identify $c\in H^1(\mathbb T^n;\R)$ with a vector in $\R^n$. Define
\[
\alpha_L(c)=-\inf_{\mu\in\mathcal I}\int(L(x,v,t)-c\cdot v)\,\dd\mu.
\]
The question is whether analyticity of $\alpha_L$ and positive definiteness of $D^2\alpha_L(c)$ for every $c\in\R^n$ force complete integrability of the Lagrangian system. The source does not specify the regularity or nonautonomous meaning of complete integrability, so a smooth Liouville conclusion and a continuous invariant-graph foliation cannot silently be identified.

The aid gives convex-duality consequences. Under the hypotheses above,
\[
\nabla\alpha_L:\R^n\to\R^n
\]
is a global analytic diffeomorphism. Its Legendre dual
\[
\beta(h)=\sup_c\{c\cdot h-\alpha_L(c)\}
\]
is analytic, $\nabla\beta=(\nabla\alpha_L)^{-1}$, and
\[
D^2\beta(h)=\bigl(D^2\alpha_L(c)\bigr)^{-1}\quad(h=\nabla\alpha_L(c)).
\]
For every minimizing measure at fixed $c$, its rotation vector $\rho(\mu)=\int v\,\dd\mu$ equals $h=\nabla\alpha_L(c)$, and its unmodified action equals $c\cdot h-\alpha_L(c)=\beta(h)$. Equality of rotation vectors and actions does not assert uniqueness of the measure or its support. Pointwise positive-definite Hessian does not add an unstated uniform lower bound on all of $\R^n$.

The aid reports a positive autonomous two-torus special case via Massart--Sorrentino: the dual regularity yields $C^0$ integrability, understood as a continuous foliation of the cotangent bundle by invariant Lagrangian graphs of the appropriate Lipschitz regularity. This is not asserted to prove smooth complete integrability for general time-periodic $n$-dimensional systems.

Its explicit spatially homogeneous time-periodic family is
\[
L_{B,a}(v,t)=\tfrac12(v-a(t))^TB(t)^{-1}(v-a(t)),
\]
where $B(t)$ is an analytic periodic symmetric positive-definite matrix and $a(t)$ an analytic periodic vector. With averages over one unit period,
\[
\overline B=\int_0^1B(t)\,\dd t,\quad \overline a=\int_0^1a(t)\,\dd t,\qquad
\alpha_{L_{B,a}}(c)=\tfrac12c^T\overline Bc+\overline a\cdot c.
\]
Momentum $p=B(t)^{-1}(v-a(t))$ is constant, and $x(t)=x(0)+\int_0^t(B(s)p+a(s))\,\dd s$ modulo $\Z^n$, giving the stated smooth integrability of this family. All coefficient pairs with the same two averages have identical $\alpha$ functions, although their time-resolved coefficients differ. This information loss is not a counterexample to the original implication, since every member of the exhibited family is integrable.
\subsection{Short English statement}
Bernard asks whether an analytic strictly convex Mather function forces complete integrability. The aid derives analytic convex duality and uniform rotation data for minimizing measures, reports a continuous-integrability theorem in the autonomous two-torus case, and gives explicitly integrable time-periodic examples whose Mather function remembers only averaged coefficients. The general target remains distinct.
\subsection{Sources}
\begin{itemize}
\item Local UnsolvedMath registry, permanent ID $20002468$, \texttt{catalogue entry}; \texttt{aim-pdes-notes.json}, \texttt{notes[79]}.
\item American Institute of Mathematics, \emph{New connections between dynamical systems and PDE's}, Question 5 (Patrick Bernard), version dated 15 August 2003 according to supplied verification: \url{https://aimath.org/WWN/dynpde/dynpde.pdf}.
\item The aid names Massart--Sorrentino for the autonomous two-torus $C^0$-integrability implication; no title, theorem number, date, or full authoritative citation is supplied.
\item Research attempt 1: partial result, problem partially solved at run, rigor label ``proved,'' novelty confidence low. Registry classification partial progress, status partially solved, statement status exact. No independent current literature assessment is supplied.
\end{itemize}
\end{document}
